\clearpage\thispagestyle{empty}
\pdfbookmark[0]{Quellen zur Besprechung von Pasquichs Sternwartenvorschlägen}{pasquich1808-sources}
\begin{center}\textbf{Quellen zur Besprechung von Pasquichs Sternwartenvorschlägen}\end{center}
\noindent Diese KI-Übertragung ins heutige Deutsch folgt der Besprechung von Gauß auf S.~497--498 der \href{https://gdz.sub.uni-goettingen.de/id/PPN236007467}{\textit{Werke}, Bd.~VI} (Göttingen, 1874). Sie erschien ohne Namensnennung unter der Überschrift \emph{Ofen} in den \textit{Göttingischen gelehrten Anzeigen}, Stück~85, 28. Mai 1808, S.~\href{https://api.digitale-sammlungen.de/iiif/image/v2/bsb10538673_00847/full/1800,/0/default.jpg}{841}--\href{https://api.digitale-sammlungen.de/iiif/image/v2/bsb10538673_00849/full/1800,/0/default.jpg}{843}. Das verlinkte Exemplar stammt aus der Bayerischen Staatsbibliothek.
\par\smallskip\noindent Gauß zitiert den Titel von Johann Pasquichs Schrift in der Form \emph{Rechenschaft von meinen Vorschlägen zur Beförderung der Astronomie auf der königl. Universitäts-Sternwarte in Ofen}. Ofen ist der deutsche Name für Buda. KI-Übertragung und KI-Satz.
\clearpage
